\subsection{The Breen comparison: scans 374--376}

We opened Breen's complete \emph{Extensions du groupe additif},
IH\'ES 48 (1978), and checked the printed statements against the
supplied transcription. All page references below are journal
pages. The manuscript scans themselves were not attached for these
leaves; consequently a correction to the supplied formula does
not decide whether the slip is in the manuscript or its AI
transcription. The main transcription locators are
\texttt{K136-03\_transcription.tex}, lines 718--945, and the inherited
reading is \texttt{136\_282-376.md}, lines 382--395.

\begin{proposition}[D013: \status{Known, used and rechecked}; hypotheses]
\label{sa08:D013}
Breen's theorems use a ring object $\mathcal O$ of characteristic
$p>0$ in a topos, with $R=\Gamma(\mathcal O)$, and hypotheses
\begin{enumerate}
\item[(A1)] maps of underlying objects $\mathcal O^d\to\mathcal O$
are precisely the polynomials $R[x_1,\ldots,x_d]$;
\item[(A2)] $H^q(\mathcal O^d,\mathcal O)=0$ for every $q>0$
and finite $d$, including $d=0$.
\end{enumerate}
Write $G_a$ for the additive underlying $\mathbb F_p$-module
sheaf. Then
\[
 H=\operatorname{End}(G_a)=R[F],\qquad Fa=a^pF,
\]
a skew polynomial ring in general. The hypotheses hold on the
big fppf, big etale and big Zariski sites over any affine
$\mathbb F_p$-algebra, and on subsites containing all affine
spaces. They do not automatically hold on the corresponding
small sites or arbitrary topoi.
\end{proposition}
\begin{proof}
These are Breen's hypotheses on pp.~45--46, Lemma~1.2 on p.~46,
and Examples~1.4--1.6 on pp.~47--48. The examples verify (A1)
by representability and Yoneda and (A2) by affine additive
cohomology vanishing; the exclusion of the small sites is explicit.
For nonaffine bases, Theorem~1.7, p.~49, gives sheaf Ext and the
local-to-global spectral sequence, not the same unrestricted
global Ext formula.
\end{proof}

We distinguish $\operatorname{Ext}_{\mathbb F_p}$ in sheaves of
$\mathbb F_p$-modules from $\operatorname{Ext}_{\mathbb Z}$ in
all abelian sheaves. Having both arguments killed by $p$ does not
identify these categories of extensions.

\begin{proposition}[D014: \status{Refuted, with corrected formulas}; coefficients]
\label{sa08:D014}
The formula on scan 374
$H^*(X;R)=H^*(X;\mathbb F_p)\otimes_{\mathbb F_p}R$ requires
finiteness hypotheses. Even when it holds, the printed pair
$P_R^i=P^i\otimes1$ for $i>0$ and $P_R^0=1\otimes F$ is not
Breen's coefficient-natural system. Its individual operations satisfy
\[
 P_R^i=P^i\otimes F,\qquad Q_R^i=Q^i\otimes F,
\]
and in characteristic two $\operatorname{Sq}_R^i=
\operatorname{Sq}^i\otimes F$, including index zero. A word of
length $\ell$ has coefficient action $F^\ell$. The identity is the
empty word and differs from $P^0=F$ on $G_a$.
\end{proposition}
\begin{proof}
For discrete $X=\mathbb N$ and $R=\mathbb F_p[t]$, the left side
in degree zero is $\prod_{n\geq0}R$. The canonical image of the
right side consists of sequences with a uniform finite polynomial
degree bound, so misses $(t^n)_n$. This is an infinite counterexample.
The tensor formula does hold if $R$ is finite dimensional over
$\mathbb F_p$, or if the relevant chain groups of $X$ are finite.

Breen gives his $p$-fold construction in (4.1)--(4.3), pp.~83--84.
The scalar relations are (4.13)--(4.14), p.~89. Alternatively
Cartan's formula applied to a degree-zero scalar $a$ gives
$P^i(ax)=a^pP^i(x)$, because the only nonzero operation on $a$
is $P^0(a)=a^p$. The same argument applies to $Q^i$ and squares.
For a finite algebraic contradiction take
$R=\mathbb F_2[t]/(t^2)$ and $X=\mathbb{CP}^2$, with
$u\in H^2(X;\mathbb F_2)$. The top-square axiom gives
\[
 \operatorname{Sq}^2(tu)=(tu)^2=0,
 \qquad t\operatorname{Sq}^2(u)=tu^2\ne0.
\]
The program checks these values exactly in
$\mathbb F_2[t,u]/(t^2,u^3)$. A separate scalar-linear extension
can be defined, but it cannot have the asserted Frobenius zeroth
operation and Cartan law simultaneously. Breen's Lemma~4.1,
p.~84, identifies $P^0$ with Frobenius, while the empty-word
convention on p.~85 identifies the unit with the identity.
\end{proof}

\begin{theorem}[D015: \status{Refuted as printed; known corrected theorem};
 cohomology operations]\label{sa08:D015}
Under (A1)--(A2), for $n\geq1$ and odd $p$,
$H^*(K(\mathcal O,n),\mathcal O)$ is the free graded-commutative
$R$-algebra on the extended admissible generators $P^J\iota_n$
with
\[
 e(J)<n\quad\text{or}\quad e(J)=n\text{ and }\epsilon_1=1.
\]
The extended operations include both $P^i$ of degree $2i(p-1)$
and $Q^i$ of degree $2i(p-1)+1$, for all $i\geq0$.
For $p=2$ the answer is the polynomial algebra on admissible
$\operatorname{Sq}^I\iota_n$ with $e(I)<n$, permitting zero
indices and retaining $\operatorname{Sq}^0$.
\end{theorem}
\begin{proof}
This is exactly Breen, Theorems~4.3--4.4, p.~86. The odd-primary
list in the supplied scan 374 omits $Q^i$ and the boundary-excess
case. The notation formally written $\beta P^i$ in an admissible
word denotes $Q^i$; it is not the composite of the actual
Bockstein with $P^i$. Breen, Remark~4.2(b), pp.~85--86, gives
\[
 \beta=Q^0,\qquad \beta P^i=Q^iP^0.
\]
His proof, pp.~86--87, starts with Lazard's calculation for $n=1$
and uses the bar/Leray spectral sequence and transgression. Thus
the manuscript's stated proof strategy is correct, with the
correct generator list. We use the opened Breen proof, not an
unopened theorem attributed to Lazard.

There is also a complete independent counterexample to the printed
odd-primary list. The normalized polynomial cochain
\[
 c_p(x,y)=\frac{(x+y)^p-x^p-y^p}{p}\pmod p
\]
is well-defined over $\mathbb F_p$. Its cocycle identity is
\[
 c_p(y,z)-c_p(x+y,z)+c_p(x,y+z)-c_p(x,y)=0,
\]
by cancellation in the integral polynomial numerator. It is not a
coboundary: the differential preserves total polynomial degree,
and the degree-$p$ part of a one-variable polynomial is a multiple
of $x^p$, whose additive coboundary is zero. But $c_p$ has the
nonzero coefficient one at $xy^{p-1}$. By (A1)--(A2) this bar
complex computes $H^*(BG_a,G_a)$, as in Breen, Lemma~2.11.
For $p=3$ it is $x^2y+xy^2$, a nonzero degree-two class.
A free algebra using only the printed even-degree $P^i$ on a
fundamental degree-one class has no such class: odd squares vanish
and positive operations increase degree by multiples of four.
The program checks the polynomial identity for $p=2,3,5,7$ as
an identity of coefficient dictionaries, not by evaluation on
$\mathbb F_p$-points. The preceding proof covers every prime.
\end{proof}

\begin{proposition}[D016: \status{Known with corrected scope};
 derived symmetric powers and operation representability]
\label{sa08:D016}
Under (A1)--(A2), Breen, Lemma~2.11, formula~(2.15), p.~61, gives
an isomorphism of graded algebras
\[
 H^*(K(\mathcal O,n),\mathcal O)
 \simeq R^*\operatorname{Sym}(R,-n).
\]
It is computed from the cosimplicial symmetric algebra on the
finite-free dual of the simplicial model. Its comparison with
left-derived $\Gamma$ is weightwise duality; a full dual of the
sum of all weights would instead give a product. Natural
cohomology operations of type $(F,n;G,q)$ are represented by
$H^q(K(F,n),G)$ in the homotopy category of simplicial sheaves.
Additivity and stability are additional requirements.
\end{proposition}
\begin{proof}
Breen, (2.13)--(2.16), p.~61, identifies the actual cochain
complex, so both the variance and the restricted-weight convention
are checked before passing to cohomology. Over $\mathbb F_p$ the
fixed-weight terms are finite dimensional, giving the stated
duality; general $R$ is obtained by its flat
$\mathbb F_p$-base change. Operation representability is Yoneda
in the homotopy category, as in Breen's section 2 and the
stable-operations discussion in section 4. A colimit over suitable
resolutions computes derived maps, rather than maps from an
arbitrary fixed nonfibrant object.

The marginal ordinary-space statement about $K(R,n)$ holds in
the finite-dimensional additive $\mathbb F_p$-vector-space case
by choosing a basis, applying the classical computation and
Kunneth, then extending coefficients. Its unrestricted infinite
version is false. Take odd $p$, $R=\mathbb F_p[t]$, and the
alternating bilinear form on its additive basis pairing
$t^{2j}$ with $t^{2j+1}$ to one for every $j$, with other pairings
zero. This infinite-rank form represents a degree-two group
cohomology class with values in $\mathbb F_p\subset R$.
After projecting coefficients by $t\mapsto0$, every finite sum
of products of degree-one $R$-valued classes has finite-rank
alternating commutator; Bockstein classes have zero commutator.
The displayed class is therefore outside the claimed generated
subalgebra. This is an infinite counterexample, certified by rank.
\end{proof}

\begin{proposition}[D017: \status{Proved}; precise finite-field detection]
\label{sa08:D017}
Over $\mathbb F_p$, the part of
$H^*(K(G_a,n),G_a)$ of total polynomial weight less than $p^h$
injects on restriction to $K(\mathbb F_{p^h},n)$ with coefficients
in $\mathbb F_{p^h}$. Each finite set of polynomial cohomology
classes is consequently detected for sufficiently large $h$.
\end{proposition}
\begin{proof}
For each classical admissible generator $I$, take its extended
lift of least polynomial weight $p^{\ell(I)}$. The generators
above it form its Frobenius string, of weights
$p^{\ell(I)+r}$, $r\geq0$, by Breen, Remark~4.2(a), p.~85,
and Remark~4.5(b), p.~87. Here $\ell(I)$ counts $p$-fold operations
in the extended lift: a $Q^i$ counts once, avoiding confusion
between the formal symbol $\beta P^i$ and a composite.
Choose an $\mathbb F_p$-basis $a_1,\ldots,a_h$ of
$\mathbb F_{p^h}$, with dual fundamental classes $z_j$.
The restriction of this extended generator is
\[
 \sum_{j=1}^h a_j^{p^{\ell(I)+r}}I_{\mathrm{cl}}(z_j),
\]
by Frobenius semilinearity and Breen's (4.8)--(4.10), pp.~87--88.
The columns of the Moore matrix
$(a_j^{p^s})_{0\leq s<h,\,1\leq j\leq h}$ are independent.
Indeed a nontrivial dependence would give an additive polynomial
of degree at most $p^{h-1}$ vanishing on all $p^h$ field elements.
This contradicts its degree. Thus all generator strings with
$\ell(I)+r<h$ inject linearly, separately for every classical
$I$. The free graded-commutative algebra on these independent
vectors embeds into the algebra on a completed basis.
D015 therefore proves the asserted weight bound. Each finite set
of classes involves finitely many generators, proving detection.
\end{proof}

The last notation on scan 374 needs this precision. Evaluation of
polynomials of degree less than $p^h$ in each variable is injective
by interpolation, but an injection of cochain groups alone does
not imply an injection in cohomology. Moreover the per-variable
bound depends on coordinates and is not preserved by arbitrary
linear coordinate changes in a simplex of $K(\mathcal O,n)$.
Thus an undefined symbol $H^*(-)_h$ cannot silently denote the
cohomology of such a non-subcomplex. D017 gives the invariant
bounded-weight statement sufficient for the detection argument.

\begin{theorem}[D018: \status{Refuted printed exponent; known correction};
 the Ext calculation]\label{sa08:D018}
Under D013, Breen's Theorem~1.3 gives
\begin{align*}
 \operatorname{Ext}^{2j+1}_{\mathbb F_p}(G_a,G_a)&=0
       &&(j\geq0),\\
 \operatorname{Ext}^{2j}_{\mathbb F_p}(G_a,G_a)
       &\simeq H/(F^{v_p(j)+1})&&(j>0).
\end{align*}
The second is an isomorphism of $(H,H^{\mathrm{op}})$-bimodules,
with the two-sided ideal generated by the displayed power.
The exponent in the supplied scan 376 is missing $+1$.
\end{theorem}
\begin{proof}
The $+1$ is explicit in Breen, Theorem~1.3, p.~47, and in the
final deduction on p.~116 from (7.22)--(7.23); both were opened,
and the first visually checked. The smallest discrepancy is
$j=1$: the supplied formula predicts $H/(1)=0$, whereas
$\operatorname{Ext}^2=H/(F)=R$, already nonzero for
$R=\mathbb F_p$. This is a counterexample by the exact opened
theorem, not an extrapolation of finite computations.

The odd vanishing includes $\operatorname{Ext}^1_{\mathbb F_p}=0$.
The inherited reading, lines 394--395, incorrectly says that this
holds only modulo $F$-torsion because of the Witt extension.
Its middle object $W_2$ is not killed by $p$:
$W_2(\mathbb F_p)=\mathbb Z/p^2$. That extension belongs to
$\operatorname{Ext}^1_{\mathbb Z}$, not
$\operatorname{Ext}^1_{\mathbb F_p}$. This resolves the apparent
contradiction completely.

The theorem computes bimodules, not the Yoneda algebra.
Breen explicitly labels the proposed multiplication
Conjecture~7.12, p.~116; the abelian-sheaf algebra is his
Conjecture~7.14, pp.~117--118. Neither algebra computation is
being inferred from the stated Ext groups.
\end{proof}

\begin{proposition}[D019: \status{Known, with categories checked};
 Frobenius torsion and perfection]\label{sa08:D019}
Every positive $\operatorname{Ext}_{\mathbb F_p}$ in D018 is
$F$-torsion. In degree $2j$ its uniform annihilator is
$F^{v_p(j)+1}$. On the perfect site over a perfect field, with
etale or Zariski topology, Breen's perfect-site corollary is
\[
 \operatorname{Ext}_{\mathbb Z}^i(G_a^{\mathrm{perf}},
              G_a^{\mathrm{perf}})=0\qquad(i>1).
\]
The attribution on scan 375 is to be refined to Breen's theorem,
with Milne as an early application and printed witness.
\end{proposition}
\begin{proof}
The torsion assertion is immediate from D018 and strengthens
$i>1$ to $i>0$ in the $\mathbb F_p$-sheaf category. The
abelian-sheaf perfect-site corollary is explicitly stated in
Breen's introduction, p.~41. Frobenius is invertible on the
perfect site, but passing Ext through perfection is justified by
that theorem, not by an unsupported formal localization argument.
Milne, \emph{Duality in the flat cohomology of a surface},
Lemma~2.2, p.~181, is an opened primary witness: it labels the
result ``L.~Breen'' and states vanishing of
$\operatorname{Ext}_{\mathbb F_p}^i(\mathcal O_S,\mathcal O_S)$
for all $i>0$ over a perfect affine base on its perfect etale
site, citing Breen's preprint and communication.
The later chapter \emph{Extensions du groupe additif sur le site
parfait}, \emph{Surfaces alg\'ebriques}, LNM 868 (1981),
pp.~238--262, was located; only publisher metadata and a two-page
preview were opened. We do not cite that unopened full chapter
as the sole evidence for the corollary.
\end{proof}

\begin{proposition}[D020: \status{Proved under the stated range};
 the spectral sequence and the two lemmas]\label{sa08:D020}
The universal-coefficient spectral sequence on scans 375--376 is
\[
 E_2^{a,b}=\operatorname{Ext}_{\mathbb F_p}^a
          (H_b(C(X)),G_a)
 \Longrightarrow H^{a+b}(X,G_a),
\]
where $C(X)$ is the free $\mathbb F_p$-linearization, with a
consistent reduced or unreduced convention. For $X=K(G_a,n)$,
$H_n=G_a$, and in the range $n\leq j<2n$ the $H_j$ are direct
sums of copies of $G_a$. The two displayed lemmas hold in this
stable range after discarding $F$-torsion. The unqualified
homology assertion in every degree is false.
\end{proposition}
\begin{proof}
Resolve the target by injectives and filter the double Hom
complex. Its two filtrations give the displayed hyper-Ext
spectral sequence. In this first-quadrant situation each total
degree involves finitely many bidegrees, so the stated use
converges; see Breen's derived Hom setup in section 2 and
fundamental spectral sequence in section 3.

For an $\mathbb F_p$-vector space $V$, Kunneth and
$(n-1)$-connectivity show that
$V\mapsto H_j(K(V,n);\mathbb F_p)$ preserves direct sums for
$j<2n$: a cross product would have degree at least $2n$.
The induced maps for sums of linear maps are therefore sums of
the induced maps. This makes the functor additive and
$\mathbb F_p$-linear, so on finite vector spaces it is
$V\otimes H_j(K(\mathbb F_p,n);\mathbb F_p)$.
Homology commutes with filtered colimits, extending the formula
to arbitrary vector spaces. Exact sheafification gives the
claimed formula for $G_a$. Breen's stabilization discussion,
p.~68, gives precisely the same strict range $j<2n$.

The unrestricted assertion fails at $n=1,j=2$.
For the finite additive groups $\mathbb F_p$ and
$\mathbb F_p^2$, the dimensions of $H_2(-;\mathbb F_p)$ are
respectively one and three, by the cyclic calculation and
Kunneth. A fixed direct sum of copies of $G_a$, whose rank is
already one on $\mathbb F_p$, would instead have dimension two
on $\mathbb F_p^2$. The supplied normalized bar matrices verify
both dimensions independently for $p=2,3$; their largest
matrix has 64 rows and 512 columns.

In the intended range, the positive Ext targets of the edge
obstruction in Lemma 1 vanish after killing $F$-torsion, by D018.
Over $\mathbb F_p$ this may be implemented by exact localization
of the commutative ring $\mathbb F_p[F]$. Over an imperfect base
we use the annihilation theorem and the stated perfect-site
result, without asserting an Ore localization.
Lemma 2 sends one $P^d$ to its classical counterpart with
coefficient $F$; an extended word of length $\ell$ contributes
$F^\ell$, including the $Q^i$. This is D014 and Breen's
restriction formulas (4.8)--(4.10). Checking $P^0$ alone, as the
manuscript suggests, would not prove the general assertion;
the $p$-fold construction and scalar relation supply that step.
\end{proof}

\begin{proposition}[D021: \status{Known/refuted according to the marked reading};
 the final Koszul calculation]\label{sa08:D021}
The extended algebra maps to the ordinary Steenrod algebra by
imposing $P^0=1$ or $\operatorname{Sq}^0=1$ and taking $Q^i$
to the classical $\beta P^i$. Its coproduct is Cartan's coproduct.
The final bar/Koszul $E_2$ computation is valid with weightwise
graded duals. Identifying that $E_2$ with the abutment for every
prime would be false.
\end{proposition}
\begin{proof}
Breen's exact weightwise Tor terms are (5.1)--(5.4), pp.~94--95;
Proposition~5.1, p.~95, computes them by a Koszul resolution.
In characteristic two all total degrees are even, so no
spectral-sequence differential survives: (5.6)--(5.7) and
Proposition~5.3, pp.~96--97. On p.~97 Breen explicitly says that
odd characteristic has nonzero differentials and $E_2$ cannot
be the abutment. Propositions~7.9 and 7.10, pp.~115--116,
compute the corrected abutment. Thus the doubtful final glyph
on scan 376 has two possible mathematical readings: a claim
only about the Koszul $E_2$ term is correct; a universal collapse
claim is false. We have not silently chosen one transcription.
Nor is a Yoneda multiplication deduced, as explained in D018.
\end{proof}

\paragraph{Corrections and scope.}
The inherited reading's odd-primary generator list (lines
383--385), its exponent $v_p(i)$ (line 390), and its interpretation
of $\operatorname{Ext}^1$ (lines 394--395) are corrected by D015
and D018. The positive-operation scalar formula at the top of
scan 374 is corrected by D014. Lead L8 distinguishes two operation
algebras but does not by itself prove that the specified
PD automorphisms realize their elements. In particular their
common unit must be distinguished from $P^0=F$; the actual
automorphism action is proved separately in Part B. All finite
computations, their scope, counts, and elapsed time are recorded
in \texttt{logs/D\_finite\_pd\_checks.json}; infinite counterexamples
are certified by the complete arguments above.
